\documentclass[10pt,a4paper]{amsart}
\usepackage[margin=19mm]{geometry}
\usepackage{amsmath,amssymb,mathtools}
\usepackage[scr=rsfs,cal=euler]{mathalfa}
\usepackage{xfrac}
\usepackage[unicode,hidelinks,bookmarks=false]{hyperref}
\newcommand{\ie}{i.e.\ }
\newcommand{\cf}{cf.\ }
\newcommand{\resp}{respectively\ }
\newcommand{\Subsection}{\subsection}
\providecommand{\nobreakdash}{\nobreak\hbox{-}}
\setlength{\emergencystretch}{2em}
\multlinegap0pt
\usepackage{bm}
\usepackage{bbm}
\usepackage{dsfont}
\usepackage{xspace}
\usepackage{cancel}
\usepackage{mathtools}
\usepackage{amsthm}
\makeatletter
\@ifundefined{proposition}{\newtheorem{proposition}{Proposition}}{}
\makeatother
\newcommand{\B}{B}
\newcommand{\Cwnp}[2]{ C_{W #1,#2}}
\newcommand{\bo}{{\rm O}}
\newcommand{\dsum}{\ds\sum}
\newcommand{\ds}{\displaystyle}
\newcommand{\eqskip}{ \vspace*{2mm}\\ }
\newcommand{\fr}[2]{\frac{\ds #1}{\ds #2}}
\newcommand{\luned}[2]{\raisebox{0.5pt}{${}$}\mathbb{L}^{#1}_{\pi/#2}}
\newcommand{\lune}[2]{\mathbb{L}^{#1}_{#2}}
\newcommand{\so}{{\rm o}}
\title[Spherical n-lunes: billiards and eigenvalues]{Spherical n-lunes: billiards and eigenvalues}
\author{Pedro Freitas, Isabel Salavessa}
\date{Source: arXiv:2608.09557v1 (CC BY 4.0)}
\begin{document}
\maketitle

\noindent\emph{Source and licence.} Spherical n-lunes: billiards and eigenvalues. Version arXiv:2608.09557v1 (CC BY 4.0), distributed under the Creative Commons Attribution 4.0 International licence, as recorded in the arXiv licence metadata for this exact version and shown on the abstract page https://arxiv.org/abs/2608.09557v1. Full rights notice: licenses/arxiv-2608.09557-CC-BY-4.0.txt.\par\smallskip

\noindent\textit{Editorial scope.} Counted unit: the Proposition whose source label is \texttt{PropC3} in the article (source0.tex lines 1854--1866, source label \texttt{PropC3}), delivered together with the article's own complete original proof of it (source0.tex lines 1867--1995, one \texttt{proof} environment of 129 lines). No mathematics is written, repaired or re-derived: statement and proof are the article's own text line for line. Required context retained uncounted: PropC.1 (lines 1698--1709). Every definition, display and label of those lines is retained unchanged. Omitted from this package and not redistributed: 1--1697, 1710--1853, 1996--3412. Nothing inside a retained excerpt is omitted or altered: every retained body is delivered exactly as the article's lines are written, so no omission receipt of any kind accompanies this package. Citations: the retained excerpts carry no citation command at all, so no bibliography entry of the article is transcribed and no reference is redistributed. Reference adaptation: none was needed and none was made; every reference inside the delivered unit resolves inside it, so no unresolved reference or citation is produced. Printed slips kept literal: no printed word, symbol, hypothesis or number of the counted statement or of its proof was altered. Layout and class: the article's own class is \texttt{amsart} with packages including amsmath, amssymb, amsfonts, eucal, amsthm, bm, bbm, geometry, array, resizegather, graphicx, tabularx, color, cancel; the wrapper is the lane's pinned \texttt{amsart} editorial wrapper at 10pt on A4, which restates the theorem-like environments the retained text needs and adds the article's own macro definitions that the retained text needs (\texttt{\textbackslash B}, \texttt{\textbackslash Cwnp}, \texttt{\textbackslash bo}, \texttt{\textbackslash ds}, \texttt{\textbackslash dsum}, \texttt{\textbackslash eqskip}, \texttt{\textbackslash fr}, \texttt{\textbackslash lune}, \texttt{\textbackslash luned}, \texttt{\textbackslash so}), with extra wrapper packages (bm, bbm, dsfont, xspace, cancel, mathtools). The delivered numbering is the unit's own. Assets: no retained span contains a figure environment, a graphics-inclusion command, a PDF-page inclusion, a table or a TikZ picture; no image file is copied into this package, the package declares no asset dependency and no image file is redistributed with it. Licence: version arXiv:2608.09557v1 of this article is distributed under the Creative Commons Attribution 4.0 International licence, as recorded in the arXiv licence metadata for this exact version and shown on the abstract page https://arxiv.org/abs/2608.09557v1. A full rights notice is delivered at licenses/arxiv-2608.09557-CC-BY-4.0.txt. Characters: every retained excerpt is pure ASCII, so the delivered engine, XeLaTeX with its default Latin Modern text font, reports no missing character.
 \begin{proposition} \label{PropC.1}
 Assume {$p\geq 1$}. The following generalised 2-term asymptotic formula holds on $\lune{2}{\pi/p}$ for $k\to +\infty$,
 \begin{equation}\label{Gen2term}
 \lambda_k=\Cwnp{2}{p} k + c(p,k)k^{1/2} +\so(k^{1/2}),
 \end{equation}
 with 
 \[c(p,k)= \sqrt{\Cwnp{2}{p}}\, p+\sqrt{\Cwnp{2}{p}}\left(1-2\frac{j}{m}\right),\]
 where $k$ is in the $K$-chain and
$j=k-k_-(K)\in \{0, \ldots, m_{2,p}(K)-1\}$. This
 second coefficient function $c(p,k)$ is bounded and  sharp along highest and lowest orders where equality is achieved, 
 \[\sqrt{\Cwnp{2}{p}}(p-1) = c(p,k_+(K))\leq c(p,k)= \sqrt{2p}\left(p+1-2\frac{j}{m}\right)\leq c(p,k_-(K))=\sqrt{\Cwnp{2}{p}}(p+1).\]
\end{proposition}

\begin{proposition} \label{PropC3} Assume $p\geq 1$ and $n\geq 3$.
The following generalised $2$-term asymptotic formula on $\luned{n}{p}$ holds for $k\to +\infty$,
\begin{eqnarray}\label{o-3term}
 \lambda_k &=&\Cwnp{n}{p}k^{2/n}+  c(p,k)k^{1/n} +\so(k^{1/n})
\end{eqnarray}
 where
 \[c(p,k)=\sqrt{\Cwnp{n}{p}}\, p+ \sqrt{\Cwnp{n}{p}}\left( 1-2\frac{j}{m_{n,p}(K)-1}\right),\]
satisfies
\begin{gather*}
\sqrt{\Cwnp{n}{p}}(p-1)=~c(p,k_+)~\leq ~ c(p,k)~\leq ~  c(p,k_-)~=~\sqrt{\Cwnp{n}{p}}(p+1).
\end{gather*}
This extends the expression of $c(p,k)$ obtained for $n=2$ in Proposition~\ref{PropC.1}
 \end{proposition}

 \begin{proof}
 %[Proof of the Proposition~\ref{PropC3}]
 %We use integers $0\leq i, \alpha\leq n-1$ and the fact that  for any family $(A_{i,\alpha})$, $i, \alpha=0, \ldots, n-1$ we have
%\[\sum_{i=0}^{n-1}\sum_{\alpha=0}^{n-1-i}A_{i,\alpha}=\sum_{0\leq i+\alpha\leq n-1}A_{i,\alpha}
%=\sum_{\alpha=0}^{n-1}\sum_{i=0}^{n-1-\alpha}A_{i,\alpha}\]
Denoting by $\sigma_i=\sigma_i(1,\ldots, n-1)$ the elementary symmetric polynomial of degree $i$ in $(n-1)$ variables,  we have
\[
%\txtg{risingelem}
(x+1)^{\overline{n-1}}=(x+1)(x+2)\ldots (x+n-1)=
\dsum_{i=0}^{n-1}\sigma_ix^{n-1-i},
\]
obtaining 
\[
\begin{array}{lll}
k_+(K) & = &  \fr{1}{(n-1)!}\dsum_{s=0}^m(sp+r+1)^{\overline{n-1}}\eqskip
& = & \dsum_{s=0}^{m}\left(  \dsum_{i=0}^{n-1}\fr{\sigma_i}{\sigma_{n-1}}(sp +r)^{n-1-i}\right) \eqskip
& = & \dsum_{s=0}^{m}\left(  \dsum_{i=0}^{n-1}\dsum_{\alpha=0}^{n-1-i}\fr{\sigma_i}{\sigma_{n-1}}\binom{n-1-i}{\alpha}(sp)^{\alpha} r^{n-1-i-\alpha}\right)\eqskip
& = & \dsum_{s=0}^{m}\left(  \dsum_{\alpha=0}^{n-1}\dsum_{i=0}^{n-1-\alpha}\frac{\sigma_i}{\sigma_{n-1}}\binom{n-1-i}{\alpha}(sp)^{\alpha} r^{n-1-i-\alpha}\right)\eqskip
& = & \dsum_{\alpha=0}^{n-1}\left(\dsum_{s=0}^{m}(sp)^{\alpha}\left[\sum_{i=0}^{n-1-\alpha}
 \frac{\sigma_i}{\sigma_{n-1}}\binom{n-1-i}{\alpha} r^{n-1-i-\alpha}\right]\right).
\end{array}
\]
Applying Faulhaber's formula we get
 \begin{eqnarray}
k_+(K) &=&  \sum_{s=0}^m(sp)^0f_0
 +\sum_{s=1}^m s f_1+
\sum_{\alpha=2}^{n-1}\left(\frac{m^{\alpha+1}}{\alpha+1}+\frac{m^{\alpha}}{2}+\frac{1}{\alpha+1}\sum_{l=2}^{\alpha}\binom{\alpha+1}{l}\B_lm^{\alpha+1-l}\right)f_{\alpha}\nonumber\eqskip
&=& (m+1)f_0+\frac{m(m+1)}{2}f_1+ \sum_{\alpha=2}^{n-1}\left(\frac{m^{\alpha+1}}{\alpha+1}+\frac{m^{\alpha}}{2}+\frac{1}{\alpha+1}
\sum_{l=2}^{\alpha}\binom{\alpha+1}{l}\B_lm^{\alpha+1-l}\right)f_{\alpha}\nonumber\eqskip
&=&
\left(\fr{m^n}{n}+\frac{m^{n-1}}{2}+\frac{1}{n}\sum_{l=2}^{n-1}\binom{n}{l}\B_lm^{n-l}\right)f_{n-1}\nonumber\eqskip
&&\hspace*{5mm}+\left(\frac{m^{n-1}}{n-1}+\frac{m^{n-2}}{2}+\frac{1}{n-1}\sum_{l=2}^{n-2}\binom{n-1}{l}\B_lm^{n-1-l}\right)f_{n-2}\nonumber\eqskip
&&\hspace*{10mm}+\ldots    \label{kmaisn}\eqskip
&&\hspace*{15mm}+\left(\frac{m^3}{3}+\frac{m^2}{2}+\frac{1}{3} \binom{3}{2}\B_2m\right)f_2+ \frac{m(m+1)}{2}f_1+(m+1)f_0\nonumber
\end{eqnarray}  
where $B_l$ are the Bernoulli numbers and    $f_{\alpha}$ are constants that  depend on $n,p,r$  given by
\[\begin{array}{ll}
f_{\alpha}=p^{\alpha}\dsum_{i=0}^{n-1-\alpha}
 \fr{\sigma_i}{\sigma_{n-1}}\binom{n-1-i}{\alpha} r^{n-1-i-\alpha}, & \alpha=0,1,2,  \ldots, n-1.\eqskip
 f_{n-1}=\fr{p^{n-1}}{(n-1)!}.
 \end{array}\]
 If $n=3$ in~\eqref{kmaisn}, then  $f_{n-1}=f_2$, $f_{n-2}=f_1$, and the corresponding coefficients coincide since the summation on $2\leq l\leq n-2=1$ is empty.
 Therefore, for any  $n\geq 3$ the following holds
\begin{eqnarray}\label{kmais2}
k_+(K) =  m^{n}\left(\frac{f_{n-1}}{n}\right) +m^{n-1}\left(\frac{f_{n-1}}{2}+\frac{f_{n-2}}{n-1}\right) +\bo(m^{n-2}),\\
\label{tilde2term}
m_{n,p}(K) = m^{n-1}\left(\frac{\tilde{f}_{n-2}}{n-1}\right) +m^{n-2}\left(\frac{\tilde{f}_{n-2}}{2}+\frac{\tilde{f}_{n-1}}{n-2}\right) +{\bo(m^{n-3})},
\end{eqnarray}
where $\bo(m^{k})$ is a polynomial of degree ${k}$ on the variable $m$, and  we use the tilde $\tilde{f}_{\alpha}$  with respect to
$m_{n,p}(K)$, and 
  $\tilde{\sigma}_i$ the $i-$elementary symmetric function on $n-2$ variables applied to $(1,2, \ldots, n-2)$. That is,  $\tilde{f}_{\beta}$ are given by
\[\begin{array}{l}
\tilde{f}_{\beta}=p^{\beta}\dsum_{i=0}^{n-2-\beta}
 \fr{\tilde{\sigma}_i}{\tilde{\sigma}_{n-2}}\binom{n-2-i}{\tilde{\beta}} r^{n-2-i-\beta}, \quad \beta=0,1, 2, \ldots, n-2.
 \end{array}\]

 Next we compute $k= k_-(K)+ j= k_+(K)-m_{n,p}(K)+1+j$, $j=0, \ldots, m_{n,p}(K)-1$.
 Observe that $m_{n,p}(K)/m^{n-1}$ is bounded. Now we have,
\begin{eqnarray}\nonumber
k^{2/n}&=&\big{(}k_+-m_{n,p}(K)+1+j \big{)}^{2/n}\eqskip
&=&\left[ m^n\left(\frac{f_{n-1}}{n}\right)+ m^{n-1}\left(\frac{f_{n-1}}{2}
+\frac{(f_{n-2}-\tilde{f}_{n-2})}{n-1} +\frac{1+j}{m^{n-1}}\right)+ \bo(m^{n-2})\right]^{2/n}\nonumber
 \eqskip \label{upper}
 &=& \left(m^n\,\frac{f_{n-1}}{n}\right)^{2/n}\left( 1 +\frac{1}{m}\left( n/2
+\frac{n(f_{n-2}-\tilde{f}_{n-2})}{(n-1)f_{n-1}} +\frac{(1+j)n}{m^{n-1}f_{n-1}}\right)+\bo\left(\frac{1}{m^2}\right)\right)^{2/n}.
 \end{eqnarray}
  Note that
 \begin{equation}\label{cut}
 (n!p)^{2/n}\left(m^n\,\frac{f_{n-1}}{n}\right)^{2/n}=m^2p^2,
 \end{equation}
 and from~\eqref{tilde2term} we obtain, % $m^{n-1}\frac{p^{n-2}}{(n-1)!}=m_{n,p}(K)+ \bo(m^{n-2})$
 \begin{equation}\label{almost}
 \frac{1}{m_{n,p}(K)}-\frac{(n-1)!}{m^{n-1}p^{n-2}}=\frac{\bo(m^{n-2})}{\bo(m^{2n-2})}=\bo\left(\frac{1}{m^{n}}\right).
 \end{equation}
Using~\eqref{upper},~\eqref{cut}) and~\eqref{almost}, and the two-term binomial expansion, we have,
\[
\begin{array}{lll}
\lambda_k-C_{W,n,p}k^{2/n} & = & Q_2(m) - m^2p^2\left[ 1 +\fr{1}{m}\left(\fr{n}{2}
+\fr{n(f_{n-2}-\tilde{f}_{n-2})}{(n-1)f_{n-1}} +\fr{(1+j)n}{m^{n-1}f_{n-1}}\right)\right.\eqskip
& &\hspace*{5mm}\left.+\bo\left(\fr{1}{m^2}\right)\right]^{2/n}\eqskip
&=& Q_2(m) - m^2p^2\left[] 1 +\fr{1}{m}\left(1+\fr{2(f_{n-2}-\tilde{f}_{n-2})}{(n-1)f_{n-1}} +\fr{2(1+j)}{m^{n-1}f_{n-1}}\right)\right.
\eqskip
& &\hspace*{5mm} \left.+\fr{2}{n}\bo(\frac{1}{m^2})+\so
\left(\fr{1}{m}\right)
\right]\eqskip
&=& mp(2p+2r+n-1)+(p+r)(p+r+n-1)\eqskip
& &\hspace*{5mm}-{mp^2}\left(1 +\fr{2(f_{n-2}-\tilde{f}_{n-2})}{(n-1)f_{n-1}} + \fr{2(1+j)}{m^{n-1}f_{n-1}}\right)+\so(m).
\end{array}
\]
Now,
 $f_{n-2}=\frac{p^{n-2}}{(n-2)!}\left(r+\frac{n}{2}\right)~$, $\frac{f_{n-2}}{f_{n-1}}=\frac{(n-1)(r+\frac{r}{2})}{p}~$, $\tilde{f}_{n-2}=\frac{p^{n-2}}{(n-2)!}$,
 and taking $k\to +\infty~$ we have
\begin{gather*}
\begin{array}{lll}
\lefteqn{\lim_{m\to +\infty}\frac{\lambda_k-C_{W,n,p}k^{2/n}}{k^{1/n}
}}\nonumber\eqskip
&=& {\ds\lim_{m\to +\infty}}\frac{ mp(2p+2r+n-1)+(p+r)(p+r+n-1)-{mp^2}\left(1
+\frac{2(f_{n-2}-\tilde{f}_{n-2})}{(n-1)f_{n-1}} +
\frac{2(1+j)}{m^{n-1}f_{n-1}}\right)+\so(m)}{
\left(m^n\,\frac{f_{n-1}}{n}\right)^{1/n}\left( 1 +\frac{1}{m}\left(\frac{1}{2}
+\frac{(f_{n-2}-\tilde{f}_{n-2})}{(n-1)f_{n-1}} +\frac{(1+j)}{m^{n-1}f_{n-1}}\right)+\bo\left(\frac{1}{m^2}\right)\right)
}\nonumber\eqskip
&=&\!\! {\ds\lim_{m\to +\infty}}
 \frac{mp(2p+2r+n-1)-{mp^2}
\left(1+2\frac{f_{n-2}}{(n-1)f_{n-1}} +
\frac{2(1+j)}{m^{n-1}f_{n-1}}
-2\frac{\tilde{f}_{n-2}}{(n-1)f_{n-1}}\right)
}
{m\left(\fr{f_{n-1}}{n}\right)^{1/n}}\eqskip
&=& {\ds\lim_{m\to +\infty}} \fr{mp(p-1)-\fr{2(1+j)(n-1)!}{m^{n-2}p^{n-3}}+2mp}{m\left(\fr{p^{n-1}}{n!}\right)^{1/n}}\eqskip
&=& {\ds\lim_{m\to +\infty}} \left[
(pn!)^{1/n}(p+1)-(pn!)^{1/n}\left(\fr{2(1+j)(n-1)!}{m^{n-1}p^{n-2}}\right)\right]\eqskip
&=&  {\ds\lim_{m\to +\infty}}
\sqrt{C_{W n,p}}\left( p+1-\fr{2(1+j)(n-1)!}{m^{n-1}p^{n-2}}\right)\eqskip
&=& {\ds\lim_{m\to +\infty}}
\sqrt{C_{W n,p}}\left( p+1-\fr{2(1+j)}{m_{n,p}(K)}+(1+ j)\times \bo\left(\fr{1}{m^n}\right)\right)  \eqskip
&=& {\ds\lim_{m \to +\infty}}
\sqrt{C_{W n,p}}\left( p+1-\fr{2j}{m_{n,p}(K)-1}\right)
\end{array}
\end{gather*}
where the second but last identity comes from~\eqref{almost},~\eqref{tilde2term}, and the fact that
 $j\leq m_{n,p}(K)-1$.
 Note that this limit may not be defined for different sublimits of $j/(m_{n,p}(K)-1)$, but if we take the following quotient we get the existence of the generalised two-term asymptotics,
\begin{eqnarray*}
\lim_{k\to +\infty}\frac{\lambda_k-C_{W,n,p}k^{2/n}- c(p,k)k^{1/n}}{k^{1/n}} &=&\lim_{k\to +\infty}\left(\frac{\lambda_k-C_{W,n,p}k^{2/n}}{k^{1/n}}- c(p,k)\right)~
 =~0,
\end{eqnarray*}
and the Proposition is proved for $n\geq 3$, extending the  same formula for the case $n=2$. 
\end{proof}

\end{document}
